\documentclass[10.5pt,a4paper]{ltjsarticle}

\usepackage[top=20mm,bottom=22mm,left=18mm,right=18mm]{geometry}
\usepackage{luatexja-fontspec}
\usepackage{fontspec}
\setmainfont{Noto Sans CJK JP}
\setsansfont{Noto Sans CJK JP}
\setmainjfont{Noto Sans CJK JP}
\setsansjfont{Noto Sans CJK JP}
\usepackage{xcolor}
\definecolor{mainblue}{HTML}{204A87}
\definecolor{lightblue}{HTML}{EAF1F8}
\definecolor{darkgray}{HTML}{333333}
\definecolor{lightgray}{HTML}{F5F5F5}
\definecolor{midgray}{HTML}{777777}

\usepackage{hyperref}
\hypersetup{colorlinks=true, linkcolor=mainblue, urlcolor=mainblue, citecolor=mainblue}
\usepackage{titlesec}
\usepackage{fancyhdr}
\usepackage{lastpage}
\usepackage{enumitem}
\usepackage{tabularx}
\usepackage{longtable}
\usepackage{array}
\usepackage{booktabs}
\usepackage[most]{tcolorbox}
\usepackage{multicol}
\usepackage{amsmath}

\setlength{\parindent}{0pt}
\setlength{\parskip}{0.55em}
\renewcommand{\baselinestretch}{1.08}
\setlist[itemize]{leftmargin=2.0em, itemsep=0.15em, topsep=0.2em}
\setlist[enumerate]{leftmargin=2.0em, itemsep=0.15em, topsep=0.2em}
\renewcommand{\arraystretch}{1.22}
\newcolumntype{Y}{>{\raggedright\arraybackslash}X}
\newcolumntype{L}[1]{>{\raggedright\arraybackslash}p{#1}}

\titleformat{\section}{\Large\bfseries\color{mainblue}}{\thesection}{0.8em}{}
\titleformat{\subsection}{\large\bfseries\color{darkgray}}{\thesubsection}{0.8em}{}
\titleformat{\subsubsection}{\normalsize\bfseries\color{darkgray}}{\thesubsubsection}{0.8em}{}
\titlespacing*{\section}{0pt}{1.1em}{0.4em}
\titlespacing*{\subsection}{0pt}{0.9em}{0.25em}
\titlespacing*{\subsubsection}{0pt}{0.6em}{0.15em}

\pagestyle{fancy}
\fancyhf{}
\lhead{03章 ソフトウェア設計}
\rhead{SWEBOK v4.0a 日本語まとめ}
\cfoot{\thepage/\pageref{LastPage}}
\renewcommand{\headrulewidth}{0.3pt}

\tcbset{
  colback=lightgray,
  colframe=midgray,
  arc=1mm,
  boxrule=0.4pt,
  left=1.0mm,
  right=1.0mm,
  top=1.0mm,
  bottom=1.0mm,
  breakable
}
\newtcolorbox{pointbox}[1]{colback=lightblue,colframe=mainblue,title=#1,fonttitle=\bfseries}
\newtcolorbox{todobox}{colback=white,colframe=midgray,title=TODO（図）,fonttitle=\bfseries}
\newtcolorbox{promptbox}{colback=lightgray,colframe=midgray,title=画像生成用プロンプト,fonttitle=\bfseries}

\newcommand{\term}[2]{\textbf{#1}\,（#2）}
\newcommand{\swebok}{\textit{SWEBOK Guide v4.0a}}

\begin{document}

\begin{titlepage}
  \vspace*{6mm}
  {\Huge\bfseries 03章 ソフトウェア設計}\par
  \vspace{3mm}
  {\Large Software Design の日本語まとめ}\par
  \vspace{3mm}
  {\normalsize \swebok{} Chapter 03 を初学者向けに再構成}\par
  \vspace{3mm}
  {\normalsize 作成日：2026年4月29日}\par

  \vspace{8mm}
  \begin{pointbox}{この資料の主張}
  ソフトウェア設計は、要求を単に「コードにする前の図」に変換する作業ではない。設計は、要求・制約・利害関係者の関心を、実装でき、評価でき、保守できる解決策へ変換する活動である。設計では、抽象化、関心事の分離、モジュール化、カプセル化、結合度、凝集度などの原則を使って複雑さを扱う。また、設計判断には根拠があり、その根拠を記録することで、後の変更や保守で同じ議論を繰り返さずに済む。
  \end{pointbox}

  \begin{pointbox}{この資料の読み方}
  専門用語は原則として日本語で示し、必要に応じて英語を括弧内に併記する。英語名は、元資料やツール上の名称を確認したいときの手がかりとして残す。初学者が前提知識なしに読めるように、各節では「何のための概念か」「設計時に何を決めるのか」「実務では何を確認するのか」を重視する。文体は、だである調で統一する。
  \end{pointbox}

  \vfill
  {\small 本資料は学習用の要約であり、原典の代替ではない。詳細な定義、参照文献、用語の範囲は原典を参照すること。}
\end{titlepage}

\tableofcontents
\newpage

\section*{全体概要：この章で学ぶこと}
\addcontentsline{toc}{section}{全体概要：この章で学ぶこと}

この章は、ソフトウェア開発における「設計」を扱う。設計は、要求を実装へ橋渡しする中心的な活動である。要求が「何を満たすべきか」を示すなら、設計は「どのような構造と振る舞いで満たすか」を示す。設計は、プログラムを書く前の一時的な準備だけではなく、保守、変更、品質評価、テスト、認証にも影響する長期的な知識資産である。

\textbf{到達目標} この章を読み終えると、次のことができるようになる。
\begin{itemize}
  \item ソフトウェア設計の意味を、設計という分野、設計活動、設計結果、ライフサイクル上の段階として説明できる。
  \item 設計思考、設計文脈、主要論点、設計原則の役割を説明できる。
  \item 高水準設計と詳細設計の違いを説明できる。
  \item 並行性、イベント処理、データ永続化、分散、例外処理、統合、安全性、変動性など、設計品質の観点を説明できる。
  \item 構造設計記述と振る舞い設計記述の代表的な記法を見分けられる。
  \item 設計パターン、領域固有言語、設計根拠、設計レビュー、設計メトリクスの意味を説明できる。
\end{itemize}

\textbf{学習順序} 初学者は、まず「設計は何を変換する活動か」を押さえる。その後、「よい設計の原則」「設計の段階」「設計で扱う品質」「設計をどう記録するか」「設計方法の違い」「設計をどう評価するか」の順に読むとよい。

\begin{todobox}
03章の全体地図を入れる。
\end{todobox}
\begin{promptbox}
白背景、モノクロ、A4本文幅に収まる横長の学習用図解を作成する。中央に「ソフトウェア設計」を置き、周囲に「基礎」「設計プロセス」「設計品質」「設計の記録」「設計戦略と方法」「品質分析と評価」を配置する。矢印は「要求から設計へ」「設計から実装・テストへ」「評価から設計改善へ」の流れを表す。ラベルは日本語、細線、講義ノート風、余白を広めにする。
\end{promptbox}

\section*{最初に必要な用語}
\addcontentsline{toc}{section}{最初に必要な用語}

\begin{tabularx}{\linewidth}{L{34mm}Y}
\toprule
\textbf{用語} & \textbf{初学者向けの説明} \\
\midrule
ソフトウェア設計 & 要求、制約、品質目標を、実装できる構造・振る舞い・インタフェースへ変換する活動である。 \\
設計記述 & 設計結果を記録した文書、モデル、図、表、プロトタイプなどである。英語では Software Design Description、略して SDD と呼ぶ。 \\
構成要素 & システムを構成する部品である。部品は、モジュール、クラス、サービス、コンポーネントなどの形を取る。 \\
インタフェース & ある構成要素が外部とやり取りする接点である。入力、出力、呼び出し方、データ形式、例外条件などを含む。 \\
抽象化 & 注目したい情報だけを取り出し、不要な詳細を隠す考え方である。複雑さを下げるために使う。 \\
モジュール化 & 大きなソフトウェアを、名前を持つ小さな部品に分けることである。理解、分担、保守を容易にする。 \\
カプセル化 & 部品の内部詳細を隠し、外部からは必要な操作だけを見せることである。情報隠蔽とも呼ぶ。 \\
品質属性 & 性能、信頼性、保守容易性、セキュリティ、使いやすさなど、ソフトウェアの性質である。 \\
トレードオフ & ある性質を高めると、別の性質や費用に悪影響が出る関係である。設計では避けられない。 \\
設計根拠 & なぜその設計判断をしたかという理由である。代替案、前提、制約、採用理由、却下理由を含む。 \\
\bottomrule
\end{tabularx}

\subsection*{略語}
\addcontentsline{toc}{subsection}{略語}

\begin{tabularx}{\linewidth}{L{18mm}L{43mm}Y}
\toprule
\textbf{略語} & \textbf{日本語} & \textbf{説明} \\
\midrule
API & アプリケーションプログラミングインタフェース & ソフトウェア部品やサービスを外部から利用するための呼び出し規約である。 \\
AOD & アスペクト指向設計 & 横断的関心事をアスペクトとして扱う設計方法である。 \\
CBD & 構成要素ベース設計 & 独立した構成要素を組み合わせてシステムを作る方法である。 \\
CRC & クラス・責務・協調カード & クラス名、責務、協調相手をカード形式で整理する方法である。 \\
DFD & データフロー図 & データの流れと処理を表す図である。 \\
DSL & 領域固有言語 & 特定領域の概念を直接表すための専用言語である。 \\
ERD & 実体関連図 & データの実体と関係を表す図である。 \\
FOSS & 自由・オープンソースソフトウェア & ソースコード利用、変更、再配布が許されるソフトウェアである。 \\
IDL & インタフェース記述言語 & 構成要素のインタフェースを記述する言語である。 \\
MBD & モデルに基づく設計 & モデルを設計記録の中心に置く方法である。 \\
MDD & モデル駆動開発 & モデルを主要成果物とし、実装やテストへつなげる開発方法である。 \\
OO & オブジェクト指向 & データと操作を持つオブジェクトを中心に考える方法である。 \\
PDL & プログラム設計言語 & プログラムに近い擬似言語で処理を表す記法である。 \\
SDD & ソフトウェア設計記述 & 設計情報を伝達・分析・実装に使うための記録である。 \\
SoC & 関心事の分離 & 複数の観点を分けて考え、複雑さを下げる原則である。 \\
UML & 統一モデリング言語 & 構造や振る舞いを表すための代表的なモデリング記法群である。 \\
\bottomrule
\end{tabularx}

\section*{導入：設計とは何か}
\addcontentsline{toc}{section}{導入：設計とは何か}

ソフトウェア設計という語は、近いが異なる意味で使われる。第一に、設計は分野である。これは、科学的原理、技術情報、想像力を使って、目的を満たすソフトウェアを経済的かつ効率的に定義する分野である。第二に、設計はその分野の中で行われるプロセスである。第三に、設計はプロセスの結果、つまり設計記述やモデルである。第四に、設計はライフサイクル上の段階でもある。

\term{ソフトウェア設計記述}{Software Design Description, SDD} は、設計情報を伝えるための成果物である。これは、実装、計画、分析、意思決定を助ける「青写真」に相当する。設計記述には、ソフトウェアを構成要素へ分けること、構成要素をどのように配置するか、構成要素同士と外部世界とのインタフェースをどのように定義するかが含まれる。

ソフトウェア設計は、要求を分析し、外部から見える性質と内部構造を定める活動である。設計は、一般に三段階で考えると理解しやすい。第一はアーキテクチャ設計であり、システム全体の根本的な構造と原則を扱う。第二は高水準設計であり、システムと主要構成要素が外部とどう関わるかを扱う。第三は詳細設計であり、各構成要素の内部を実装できる水準まで具体化する。

\begin{pointbox}{重要}
設計は「実装前にきれいな図を描くこと」ではない。設計は、要求、制約、品質目標、技術選択、保守の見通しを踏まえ、作れる解決策を定義する活動である。
\end{pointbox}

\section{ソフトウェア設計の基礎}

ソフトウェア設計の基礎は、設計を理解するための概念、用語、判断基準である。初学者にとって重要なのは、設計を「コードを書く前の一工程」と狭く見るのではなく、「問題を解決可能な形へ変換する思考」として見ることである。

\subsection{設計思考}

\term{設計思考}{Design Thinking} は、必要や問題を理解し、それに対する解決策を考える思考である。設計対象はソフトウェアに限らない。建築、工業製品、組織、業務手順など、人が目的を持って作るものには設計が関わる。

ソフトウェア設計で重要なのは、問題と解決策を表す語彙を作ることである。ソフトウェアは、物理的な材料を削って作るものではなく、概念、名前、規則、手続き、データ、インタフェースを組み合わせて作る。したがって、設計は言語的な活動でもある。

設計思考は、次の五段階として理解できる。
\begin{enumerate}
  \item 目的や目標を明確にする。
  \item その目的を達成する概念的な構造を考える。
  \item 概念構造を実現する仕組みを考える。
  \item その仕組みを表し、利用するための記法や語彙を作る。
  \item 具体的な問題文脈で記法を使い、目的を達成する。
\end{enumerate}

\begin{pointbox}{定義}
設計思考とは、問題を理解し、制約の中で複数の解決策を考え、根拠を持って選び、表現可能な形へまとめる思考である。
\end{pointbox}

\begin{todobox}
設計思考の五段階を示す図を入れる。
\end{todobox}
\begin{promptbox}
白背景、モノクロ、A4本文幅の横長フロー図を作成する。「目的の明確化」「概念構造」「仕組み」「記法・語彙」「具体的適用」の五つの箱を左から右へ並べる。上に「問題を理解する」、下に「解決策を実装可能にする」という注記を入れる。各箱に短い説明を添える。日本語ラベル、細線、講義ノート風。
\end{promptbox}

\subsection{ソフトウェア設計の文脈}

ソフトウェア設計は、開発ライフサイクルの中で孤立して存在しない。要求、アーキテクチャ、実装、テスト、保守と強くつながっている。

\begin{itemize}
  \item \textbf{要求との関係}：要求は、設計が解くべき問題を示す。設計は、要求を実装可能な構造と振る舞いへ変換する。
  \item \textbf{アーキテクチャとの関係}：アーキテクチャが定められている場合、それは設計の制約となる。主要構成要素、接続、API、採用すべきスタイルやパターン、設計原則などを制約する。
  \item \textbf{実装との関係}：設計は、実装者が何を作るべきかを示す案内である。実装者がコードを読む前に、部品の責務やインタフェースを理解できることが望ましい。
  \item \textbf{テストとの関係}：設計は、テスト戦略とテストケースの土台になる。設計要素が要求をどう満たすかが明らかであれば、テストの観点も作りやすい。
\end{itemize}

設計の形式性は、プロジェクトによって異なる。安全性が重要なシステムでは、設計記述を厳密に残す必要がある。一方、小規模な探索的開発では、簡潔なモデルやコードそのものが設計記録の中心になる場合もある。ただし、どの場合でも、関係者が設計判断を理解できることが重要である。

\begin{todobox}
設計と他知識領域の関係図を入れる。
\end{todobox}
\begin{promptbox}
白背景、モノクロ、A4本文幅の概念図を作成する。中央に「ソフトウェア設計」を置く。左に「要求」、上に「アーキテクチャ」、右に「実装」、下に「テスト」、右下に「保守」を配置する。要求から設計へ、アーキテクチャから設計へ、設計から実装・テストへ、保守から設計改善へ矢印を描く。各矢印に「問題」「制約」「青写真」「確認観点」「変更知識」と短く注記する。
\end{promptbox}

\subsection{ソフトウェア設計の主要論点}

設計では、多くの論点を同時に扱う必要がある。代表的な論点は、性能、セキュリティ、信頼性、使いやすさ、保守容易性などの品質である。また、ソフトウェアをどのような構成要素へ分けるか、それらをどう接続するか、どの単位で責務を持たせるかも重要である。

設計論点には、機能そのものではないが、システム全体に影響するものもある。たとえば、ログ、認証、監査、例外処理、キャッシュ、トランザクション、セキュリティ確認などは、多くの機能に横断的に関わる。こうした論点は\term{横断的関心事}{crosscutting concern} と呼ばれる。

\begin{pointbox}{注意}
横断的関心事を場当たり的に実装すると、同じ処理が各所に重複し、変更時に漏れが起きやすくなる。設計段階で共通方針を決めることが重要である。
\end{pointbox}

\subsection{ソフトウェア設計原則}

設計原則は、設計中の判断に方向を与える基本的な考え方である。原則は、特定のプログラミング言語や開発手法だけに依存しない。むしろ、複雑な問題を扱うための共通の知恵である。

\begin{longtable}{L{34mm}L{55mm}L{76mm}}
\toprule
\textbf{原則} & \textbf{意味} & \textbf{実務で見ること} \\
\midrule
\endhead
抽象化 & 目的に関係する情報だけを示し、不要な詳細を隠す。 & 利用者、設計者、実装者が必要な詳細水準だけを見られるか。 \\
関心事の分離 & 異なる関心を分けて扱う。 & 業務処理、認証、ログ、表示、保存などが混ざりすぎていないか。 \\
モジュール化 & 大きなソフトウェアを小さな部品へ分ける。 & 各部品の責務が明確で、独立して理解・変更できるか。 \\
カプセル化・情報隠蔽 & 内部詳細を隠し、必要なインタフェースだけを見せる。 & 外部が内部データ構造や実装順序に依存していないか。 \\
インタフェースと実装の分離 & 外部に見せる契約と内部実装を分ける。 & 実装を変えても、利用側の変更が少なく済むか。 \\
結合度 & 部品同士の依存の強さである。 & ひとつの部品変更が多数の部品変更を引き起こさないか。 \\
凝集度 & ひとつの部品内の要素がどれだけ関連しているかである。 & ひとつの部品が複数の無関係な責務を抱えていないか。 \\
一貫性 & 名前、記法、順序、インタフェースなどをそろえる。 & 似た問題に似た解決策を使っているか。 \\
完全性 & 重要な特徴や要求を漏らさない。 & 通常時、例外時、境界条件、状態、モードが考慮されているか。 \\
検証可能性 & 要求や制約を満たすか確認できる。 & 設計からテスト、レビュー、分析が可能か。 \\
倫理的配慮 & 人権、福祉、透明性、説明責任、悪用可能性などを考慮する。 & AI、自律システム、社会的影響があるシステムで設計上の配慮があるか。 \\
\bottomrule
\end{longtable}

特に、結合度と凝集度は初学者が最初に押さえるべき設計観点である。一般に、よい設計は、部品間の結合度を低くし、部品内の凝集度を高くする。結合度が高いと変更の影響が広がる。凝集度が低いと、部品が何を担当しているのか理解しにくくなる。

\begin{pointbox}{重要}
設計原則は、単独で使うものではない。抽象化、関心事の分離、モジュール化、カプセル化、低結合、高凝集は互いに支え合う。
\end{pointbox}

\begin{todobox}
主要な設計原則の関係図を入れる。
\end{todobox}
\begin{promptbox}
白背景、モノクロ、A4本文幅の関係図を作成する。中央に「複雑さを管理する」を置く。周囲に「抽象化」「関心事の分離」「モジュール化」「カプセル化」「低結合」「高凝集」「検証可能性」を配置する。抽象化からモジュール化、モジュール化から低結合・高凝集へ矢印を描く。日本語ラベル、細線、講義ノート風。
\end{promptbox}

\section{ソフトウェア設計プロセス}

ソフトウェア設計は、一般に多段階の活動として考えられる。段階は、常に厳密に分かれるわけではないが、何を決める段階かを理解すると、設計成果物を作りやすくなる。

設計は、三つの段階に分けられる。
\begin{itemize}
  \item \textbf{アーキテクチャ設計}：システム全体の根本的な構造、主要構成要素、接続方式、支配的な品質方針を決める。
  \item \textbf{高水準設計}：外部に向いた設計であり、主要構成要素と外部環境のやり取り、構成要素間のやり取りを決める。
  \item \textbf{詳細設計}：内部に向いた設計であり、構成要素の内部構造、データ構造、アルゴリズム、内部モジュールを決める。
\end{itemize}

段階の境界は、プロジェクトやシステムの性質で変わる。たとえば、通常ならプログラマが選ぶソートアルゴリズムでも、極端な性能要求があるシステムでは、初期の設計段階で決める必要がある。このように、何をいつ決めるかは、その判断が要求達成にどれほど重要かで決まる。

\begin{todobox}
設計の三段階を示す図を入れる。
\end{todobox}
\begin{promptbox}
白背景、モノクロ、A4本文幅の縦長図を作成する。上から「要求・制約」「アーキテクチャ設計」「高水準設計」「詳細設計」「実装」の箱を並べる。各段階の右に「全体方針」「外部とのやり取り」「内部構造」「コード化」と注記する。下に行くほど詳細度が増すことを左側の矢印で示す。日本語ラベル、講義ノート風。
\end{promptbox}

\subsection{高水準設計}

\term{高水準設計}{High-Level Design} は、システムの主要構成要素が互いにどう関わり、外部環境とどうやり取りするかを決める設計である。外部環境には、利用者、外部システム、デバイス、ネットワーク、データベース、運用環境などが含まれる。

高水準設計では、主に次を扱う。
\begin{itemize}
  \item システムが反応すべき外部イベントやメッセージ。
  \item システムが生成すべきイベントやメッセージ。
  \item イベントやメッセージのデータ形式とプロトコル。
  \item 入力と出力の順序、タイミング、依存関係。
  \item 端から端までの取引やイベント列の流れ。
  \item データをどのように保存し、管理するか。
\end{itemize}

高水準設計は、アーキテクチャがある場合、その枠内で行われる。たとえば、アーキテクチャが「非同期メッセージでやり取りする」と決めていれば、高水準設計ではその方針に従って、具体的なメッセージ種類、順序、例外時の扱いを決める。

\subsection{詳細設計}

\term{詳細設計}{Detailed Design} は、主要構成要素の内部を実装できる水準まで具体化する設計である。これは、プログラマがコードを書くための直接的な手がかりになる。

詳細設計では、主に次を扱う。
\begin{itemize}
  \item 主要構成要素を内部モジュールやプログラム単位へ分ける。
  \item 各モジュールの責務を割り当てる。
  \item モジュール間のやり取りを決める。
  \item 名前の有効範囲、可視性、公開・非公開の境界を決める。
  \item 状態、モード、状態遷移を定義する。
  \item データと制御の依存関係を明らかにする。
  \item データ構造、アルゴリズム、ユーザインタフェースの詳細を決める。
\end{itemize}

詳細設計は、実装と完全に分離して行われるとは限らない。特に反復型やアジャイル型の開発では、設計と実装が近い間隔で進む。しかし、その場合でも、設計判断が存在しなくなるわけではない。設計判断がコードの中だけに閉じ込められるのか、明示的に記録されるのかが変わるだけである。

\section{ソフトウェア設計品質}

設計品質とは、設計がどのような性質を持つか、またその設計によって作られるソフトウェアがどのような品質を達成しやすいかに関する観点である。ここで扱う品質は、単に「よく動く」ことではなく、並行処理、イベント処理、永続化、分散、例外処理、統合、安全性、変動性など、多面的である。

\subsection{並行性}

\term{並行性}{Concurrency} は、複数の処理が同時または重なり合って進む性質である。設計では、処理をプロセス、タスク、スレッドなどの単位へどう分けるかを考える。

並行性の設計では、効率だけでなく、同期、排他制御、原子性、スケジューリングを考慮する必要がある。たとえば、複数の処理が同じデータを書き換える場合、順序が不適切だと結果が不安定になる。したがって、どのデータを共有するか、どこでロックするか、どの処理を非同期にするかを設計で明確にする。

\subsection{制御とイベント処理}

\term{制御}{Control} は、処理の流れをどう組み立てるかに関する設計である。\term{イベント処理}{Event Handling} は、利用者操作、外部メッセージ、タイマー、センサー入力などの出来事にどう反応するかに関する設計である。

イベント処理では、同期呼び出し、非同期メッセージ、コールバック、暗黙呼び出しなどの仕組みが使われる。反応型のシステムでは、イベント発生時の順序、重複、遅延、失敗時の扱いを設計しなければならない。

\subsection{データ永続化}

\term{データ永続化}{Data Persistence} は、プログラム終了後もデータを保持し、必要なときに取り出せるようにすることである。設計では、データベース、ファイル、キャッシュ、メッセージキュー、外部ストレージなどをどう使うかを決める。

永続化の設計では、保存形式、整合性、トランザクション、検索性能、バックアップ、復旧、権限管理を考慮する。データは長期に残るため、後から変更しにくい。したがって、データ構造は実装の都合だけでなく、将来の変更と保守を見越して決める必要がある。

\subsection{構成要素の分散}

\term{構成要素の分散}{Distribution of Components} は、ソフトウェア構成要素を複数のコンピュータ、ネットワーク、デバイスへ配置する設計である。分散させると、拡張性や可用性を高められる場合があるが、通信遅延、障害、整合性、監視、復旧が難しくなる。

分散設計では、構成要素間の通信方式、データ複製、障害時の再試行、サービス発見、負荷分散、セキュリティ境界を考える。特に、ネットワーク越しの呼び出しは、ローカル関数呼び出しとは違い、遅延や失敗が普通に起こる前提で設計する必要がある。

\subsection{誤り・例外処理と耐故障性}

\term{誤り処理}{Error Handling} と\term{例外処理}{Exception Handling} は、通常の処理から外れた状況をどう扱うかに関する設計である。\term{耐故障性}{Fault Tolerance} は、故障や誤りが起きても、影響を抑え、必要なサービスを続ける性質である。

設計では、どの異常を防ぐか、どの異常を検出するか、どの異常を許容するか、どの異常では停止するかを決める。例外を単に捕捉して無視すると、原因が隠れ、後で大きな障害になる。したがって、記録、通知、復旧、再試行、代替処理、停止条件を設計で明確にする。

\subsection{統合と相互運用性}

\term{統合}{Integration} は、複数の部品やシステムを組み合わせて一つの目的を達成することである。\term{相互運用性}{Interoperability} は、異なるシステムや構成要素がデータやサービスをやり取りして協調できる性質である。

この論点は、企業内システム、システムオブシステムズ、クラウドサービス連携、異なるフレームワークやライブラリを使う場合に重要である。設計では、データ形式、プロトコル、API、認証方式、エラー通知、バージョン互換性を決める。

\subsection{保証・セキュリティ・安全性}

\term{保証}{Assurance} は、ソフトウェアが意図どおりに動くという信頼を得ることに関わる。\term{セキュリティ}{Security} は、情報や資源を不正なアクセス、変更、破壊、利用不能化から守ることである。\term{安全性}{Safety} は、人命、財産、環境への危害を避けることである。

セキュリティ設計では、機密性、完全性、可用性、認証、認可、監査、攻撃時の被害限定、復旧を考える。安全性設計では、危険状態に至る条件、停止方法、警告、フェイルセーフ、冗長化を考える。どちらも後付けが難しいため、設計段階から扱う必要がある。

\subsection{変動性}

\term{変動性}{Variability} は、ソフトウェアに許されるバリエーションを扱う設計観点である。市場、利用組織、地域、法制度、設定、利用文脈によって、同じソフトウェアに異なる振る舞いや構成が求められることがある。

変動性は、ソフトウェア製品系列やシステムファミリで特に重要である。共通部分と可変部分を分ければ、複数の派生製品を効率よく作れる。フィーチャモデルは、機能や選択肢、その依存関係を整理するために使われる。

\begin{todobox}
設計品質の関心事マップを入れる。
\end{todobox}
\begin{promptbox}
白背景、モノクロ、A4本文幅のマインドマップを作成する。中央に「設計品質」を置き、枝として「並行性」「イベント処理」「永続化」「分散」「例外・耐故障」「統合・相互運用」「保証・セキュリティ・安全」「変動性」を配置する。各枝に一言の注意点を添える。日本語、講義ノート風、細線。
\end{promptbox}

\section{ソフトウェア設計の記録}

設計プロセスの成果は、設計者の頭の中だけに残しておくべきではない。設計記録は、問題、解決策、判断、根拠を関係者へ伝えるためのものである。設計記録には、文章、図、表、モデル、プロトタイプ、コード片などが含まれる。

設計記録の目的は、実装者に作るべきものを示すことだけではない。顧客、テスト担当、品質保証担当、認証担当、保守担当も設計情報を必要とする。したがって、どの設計情報を、誰に向けて、どの詳細度で記録するかを意識して決める必要がある。

\begin{pointbox}{基本方針}
設計記録は「すべてを詳細に書く」ことが目的ではない。読者、利用目的、リスク、保守期間に応じて、必要な設計知識を取り出しやすい形で残すことが目的である。
\end{pointbox}

\subsection{モデルに基づく設計}

\term{モデルに基づく設計}{Model-Based Design, MBD} は、モデルを設計記録の中心に置く方法である。従来の文書中心の設計では、自然言語と非形式的な図を使うため、曖昧さや不完全さが残りやすい。また、必要な情報が複数文書に分散し、理解や分析が難しくなる。

モデルに基づく設計では、ツールが設計情報を整理し、表示し、分析できる。たとえば、モデルの一部をアニメーションやシミュレーションで確認したり、仮定を変えた場合の影響を調べたり、要求から設計、設計からテストへの追跡を支援したりできる。

\term{モデル駆動開発}{Model-Driven Development, MDD} は、モデルを開発プロセスの主要成果物として使い、コード生成や検証へつなげる考え方である。

\begin{todobox}
文書中心設計とモデル中心設計の比較図を入れる。
\end{todobox}
\begin{promptbox}
白背景、モノクロ、A4本文幅の左右比較図を作成する。左に「文書中心設計」として「文章」「図」「表」が分散している様子を描く。右に「モデルに基づく設計」として中央の「モデル」から「ビュー」「分析」「シミュレーション」「追跡」へ矢印が出る様子を描く。日本語ラベル、講義ノート風、細線。
\end{promptbox}

\subsection{構造設計記述}

\term{構造設計記述}{Structural Design Description} は、ソフトウェアの静的な構造を表す。つまり、どの構成要素があり、それらがどう関係し、どの責務を持つかを示す。

代表的な記法は次のとおりである。

\begin{longtable}{L{42mm}L{112mm}}
\toprule
\textbf{記法} & \textbf{説明} \\
\midrule
\endhead
クラス図・オブジェクト図 & クラス、オブジェクト、属性、操作、関連を表す。オブジェクト指向設計でよく使われる。 \\
コンポーネント図 & 交換可能な構成要素と、その提供・要求インタフェースを表す。 \\
CRCカード & クラス名、責務、協調相手をカード形式で整理する。初期設計や会話に向く。 \\
配置図 & ソフトウェアをどの物理ノードや実行環境へ配置するかを表す。 \\
実体関連図 & データの実体、属性、関係を表す。データ設計でよく使われる。 \\
インタフェース記述言語 & 構成要素が公開する操作名、型、引数、戻り値などを記述する。 \\
構造チャート & プログラムの呼び出し構造を表す。どのモジュールがどのモジュールを呼ぶかを示す。 \\
\bottomrule
\end{longtable}

\subsection{振る舞い設計記述}

\term{振る舞い設計記述}{Behavioral Design Description} は、ソフトウェアが時間の中でどのように動くかを表す。処理の流れ、イベントへの反応、状態変化、入力と出力の関係などを示す。

代表的な記法は次のとおりである。

\begin{longtable}{L{42mm}L{112mm}}
\toprule
\textbf{記法} & \textbf{説明} \\
\midrule
\endhead
活動図 & 処理の流れ、分岐、並行活動、入力、出力を表す。 \\
相互作用図 & オブジェクトや構成要素間のメッセージ交換を表す。通信図とシーケンス図が代表である。 \\
データフロー図 & 処理要素間を流れるデータを表す。セキュリティ分析にも使える。 \\
決定表・決定図 & 条件の組合せと、それに対応する動作を表す。複雑な業務規則に向く。 \\
フローチャート & 制御の流れと処理順序を表す。 \\
状態遷移図・状態チャート & 状態、イベント、遷移、状態ごとの振る舞いを表す。 \\
形式仕様言語 & 数学的な記法で、インタフェースや振る舞いを厳密に定義する。 \\
擬似コード・プログラム設計言語 & プログラムに近い形で処理手順を表す。詳細設計やアルゴリズム説明に使う。 \\
\bottomrule
\end{longtable}

\begin{todobox}
構造設計記述と振る舞い設計記述の分類図を入れる。
\end{todobox}
\begin{promptbox}
白背景、モノクロ、A4本文幅の二分割図を作成する。左に「構造を表す記法」として「クラス図」「コンポーネント図」「配置図」「ERD」「IDL」を並べる。右に「振る舞いを表す記法」として「活動図」「シーケンス図」「DFD」「決定表」「状態遷移図」「擬似コード」を並べる。中央に「設計記述」を置き、左右へ分岐する。日本語、講義ノート風。
\end{promptbox}

\subsection{設計パターンとスタイル}

\term{設計パターン}{Design Pattern} は、特定の文脈で繰り返し現れる問題に対する一般的な解決策である。設計パターンは、過去に有効だった設計上の慣用句を共有するための語彙である。

代表的な分類は、生成、構造、振る舞いである。
\begin{itemize}
  \item \textbf{生成パターン}：オブジェクトや部品の生成方法を扱う。ビルダー、ファクトリ、プロトタイプ、シングルトンなどがある。
  \item \textbf{構造パターン}：部品の組み合わせ方を扱う。アダプタ、ブリッジ、コンポジット、デコレータ、ファサード、プロキシなどがある。
  \item \textbf{振る舞いパターン}：責務の分担や相互作用を扱う。コマンド、イテレータ、メディエータ、オブザーバ、状態、戦略、テンプレート、ビジタなどがある。
\end{itemize}

\term{スタイル}{Style} は、より大きな構造の特徴を表す場合に使われることが多い。たとえば、階層化、パイプ・アンド・フィルタ、クライアントサーバ、イベント駆動などは、システム全体または大きな部分の設計方針として現れる。

\begin{todobox}
設計パターン分類の図を入れる。
\end{todobox}
\begin{promptbox}
白背景、モノクロ、A4本文幅のツリー図を作成する。最上位に「設計パターン」。下位に「生成」「構造」「振る舞い」。生成には「ファクトリ」「ビルダー」「シングルトン」、構造には「アダプタ」「ファサード」「プロキシ」、振る舞いには「オブザーバ」「戦略」「状態」を例として配置する。日本語ラベル、講義ノート風。
\end{promptbox}

\subsection{専用言語と領域固有言語}

設計表現は、構造と振る舞いの二分法だけで整理できるとは限らない。たとえば、利用者インタフェース設計では、画面の構造、操作の流れ、入力制約、表示規則が混ざる。また、安全性や信頼性の領域では、その分野に特化した記法が発達している。

\term{領域固有言語}{Domain-Specific Language, DSL} は、特定の業務領域や技術領域の概念を直接表すための言語である。汎用プログラミング言語ではなく、その領域の語彙を使って設計や実装に近い表現を行う。例として、シミュレーション、リアルタイムシステム、ゲーム開発、利用者インタフェース、テスト開発、言語処理などの領域向けDSLがある。

DSLの利点は、領域の専門家が理解しやすい語彙で表現できることである。一方、DSLを作るには、言語仕様、構文解析、生成、保守の負担がある。したがって、繰り返し使う価値がある領域で特に有効である。

\subsection{設計根拠}

\term{設計根拠}{Design Rationale} は、なぜその設計判断をしたかを説明する情報である。設計根拠には、前提、検討した代替案、採用基準、トレードオフ、却下した案の理由が含まれる。

設計チームにとって、現在の判断理由は明らかに見えることが多い。しかし、数か月後、数年後、または担当者交代後には、その理由が失われる。設計根拠を記録すれば、保守担当者は「なぜこの構造なのか」「なぜ別案を採らなかったのか」を理解しやすい。

却下した判断を記録することも有用である。前提や制約が変わったとき、以前は却下された案が有効になる場合がある。逆に、過去に明確な理由で却下された案を、理由を忘れて再び採用してしまうことを防げる。

\begin{pointbox}{実務の見方}
設計根拠は、巨大な説明文である必要はない。「判断」「理由」「代替案」「影響」「未解決事項」を短く残すだけでも、後の保守で大きな価値を持つ。
\end{pointbox}

\section{ソフトウェア設計戦略と方法}

設計戦略と方法は、設計プロセスを進めるための考え方である。方法によって、何を最初に見るか、どのように分解するか、どの概念を中心に置くかが異なる。

\subsection{一般戦略}

一般的な設計戦略には、分割統治、段階的詳細化、トップダウン、ボトムアップ、反復、増分、パターン利用などがある。

\term{分割統治}{divide and conquer} は、大きな問題を小さな問題に分けて解く考え方である。\term{段階的詳細化}{stepwise refinement} は、粗い設計から始め、徐々に詳細を追加する考え方である。トップダウンは全体から部分へ進む。ボトムアップは既存部品や低水準機能から全体を組み立てる。

実務では、これらは単独で使われるというより、組み合わせて使われる。最初にトップダウンで全体構造を作り、既存部品をボトムアップで組み込み、反復しながら詳細化することが多い。

\subsection{機能指向・構造化設計}

\term{機能指向設計}{Function-Oriented Design} または\term{構造化設計}{Structured Design} は、ソフトウェアの主要機能を見つけ、それを上位から下位へ分解する古典的な方法である。

この方法では、データフロー図や処理記述を使い、入力がどの処理を通って出力へ変換されるかを考える。業務処理の流れが明確で、機能分解しやすいシステムで理解しやすい。一方、状態やデータ構造、イベント、オブジェクト間の相互作用が複雑な場合は、別の方法と組み合わせる必要がある。

\subsection{データ中心設計}

\term{データ中心設計}{Data-Centered Design} は、プログラムが扱うデータ構造から設計を始める方法である。入力データ構造と出力データ構造を明らかにし、入力を出力へ変換する処理単位を設計する。

データが安定しており、処理がデータ構造に強く依存する場合に有効である。たとえば、ファイル変換、帳票生成、データ処理パイプラインなどでは、データ構造を中心に考えると設計しやすい。

\subsection{オブジェクト指向設計}

\term{オブジェクト指向設計}{Object-Oriented Design} は、オブジェクトを中心に設計する方法である。オブジェクトは、データと操作を一体として持つ概念である。設計では、クラス、属性、メソッド、関連、継承、多態性、責務を考える。

初期のオブジェクト指向設計では、名詞からオブジェクトを見つけ、動詞からメソッドを見つける方法が使われた。現在では、単に名詞を拾うだけでなく、責務を中心に設計する考え方が重要である。クラスが何を知り、何を行い、誰と協調するかを考える。

SOLID原則は、クラス設計でよく参照される考え方である。単一責任、開放閉鎖、リスコフ置換、インタフェース分離、依存関係逆転の頭文字である。これらは、変更しやすく理解しやすいオブジェクト指向設計を目指すための指針である。

\subsection{利用者中心設計}

\term{利用者中心設計}{User-Centered Design} は、利用者とその作業、目的、文脈を深く理解したうえで設計する方法である。単に画面をきれいにすることではなく、利用者が何を達成したいか、どのような環境で使うか、どのような誤りを起こしやすいかを考える。

利用者中心設計では、利用者要求の収集、作業の流れの整理、プロトタイプやモックアップの作成、利用者評価を行う。使いやすさが重要なシステムでは、利用者中心設計を後から追加するのではなく、設計の初期から組み込む必要がある。

\subsection{構成要素ベース設計}

\term{構成要素ベース設計}{Component-Based Design, CBD} は、独立して配置・利用できる構成要素を組み合わせてシステムを作る方法である。各構成要素は、明確なインタフェースと依存関係を持つ。

この方法の目的は、再利用、独立開発、交換容易性、統合容易性を高めることである。構成要素ベース設計では、共通API、サービスごとのAPI、部品のバージョン、依存関係、配置方法が重要になる。

\subsection{イベント駆動設計}

\term{イベント駆動設計}{Event-Driven Design} は、イベントに反応して処理を行う設計方法である。イベントとは、利用者操作、メッセージ受信、状態変化、タイマー発火、センサー入力などである。

イベント駆動設計では、発行・購読方式がよく使われる。送信側はイベントを発行し、受信側は関心のあるイベントを購読する。これにより、送信側と受信側の結合を弱められる。大規模で拡張性が必要なシステムでは有効だが、イベントの順序、重複、遅延、失敗時の影響を慎重に設計する必要がある。

\subsection{アスペクト指向設計}

\term{アスペクト指向設計}{Aspect-Oriented Design, AOD} は、横断的関心事をアスペクトとして分離する方法である。横断的関心事とは、ログ、セキュリティ確認、トランザクション、例外処理など、複数の機能にまたがる関心である。

アスペクト指向の考え方は、必ずしも独立した設計手法として広く使われているわけではない。しかし、アプリケーションフレームワークやライブラリでは、宣言や設定によって横断的な振る舞いを追加する形で利用されることがある。

\subsection{制約に基づく設計}

\term{制約に基づく設計}{Constraint-Based Design} は、制約を使って設計空間を狭める方法である。設計空間とは、可能な設計案全体の集合である。制約は、不可能または受け入れられない案を除外する役割を持つ。

制約には、ハードウェア、ソフトウェア、データ、操作手順、インタフェース、法令、性能、費用、納期などが含まれる。制約を明確にすれば、設計探索が速くなる。一方、制約が多すぎる場合や互いに矛盾する場合は、解が見つからないこともある。

\subsection{ドメイン駆動設計}

\term{ドメイン駆動設計}{Domain-Driven Design} は、対象領域の概念と言葉を中心に設計する方法である。ここでいうドメインとは、銀行、医療、物流、教育、ゲームなど、ソフトウェアが扱う問題領域である。

ドメイン駆動設計では、分析者、設計者、開発者、業務専門家が共有する言葉を重視する。この共有語彙を使って、オブジェクト、役割、イベント、活動、規則を設計記述に反映する。言葉がずれると、実装もずれるため、ドメインの語彙を設計の中心に置くことが重要である。

\subsection{その他の方法}

その他にも、多くの設計方法がある。反復型・適応型の方法では、厳密な要求と設計を最初にすべて固定するのではなく、ソフトウェア増分を作りながら設計を調整する。

\term{サービス指向設計}{Service-Oriented Design} は、分散されたサービスを組み合わせてソフトウェアを作る方法である。HTTP、HTTPS、SOAPなどの標準プロトコルを使い、異なる提供者のサービスをつなぐことがある。

\begin{todobox}
設計戦略と中心概念の対応図を入れる。
\end{todobox}
\begin{promptbox}
白背景、モノクロ、A4本文幅のマトリクス図を作成する。左列に「機能指向」「データ中心」「オブジェクト指向」「利用者中心」「構成要素ベース」「イベント駆動」「アスペクト指向」「ドメイン駆動」。右列に「中心に置くもの」を置き、それぞれ「機能」「データ構造」「責務を持つオブジェクト」「利用者の作業」「部品とインタフェース」「イベント」「横断的関心事」「領域の語彙」と対応させる。日本語、講義ノート風。
\end{promptbox}

\section{ソフトウェア設計品質分析と評価}

設計は、作ったら終わりではない。設計が要求を満たすか、品質属性を達成できるか、保守しやすいか、リスクが残っていないかを分析・評価する必要がある。

\subsection{設計レビューと監査}

\term{設計レビュー}{Design Review} は、設計を包括的に調べる活動である。目的は、設計の完成度、要求の網羅、未解決事項、潜在問題、品質リスクを確認することである。レビューは、設計チーム自身、独立した第三者、顧客、品質保証担当などによって行われる。

\term{設計監査}{Design Audit} は、より狭い観点に焦点を当てる。たとえば、特定の機能監査、標準への適合監査、セキュリティ方針への適合監査などである。

\begin{todobox}
設計レビューの流れを入れる。
\end{todobox}
\begin{promptbox}
白背景、モノクロ、A4本文幅のフローチャートを作成する。「レビュー対象の準備」「観点・チェックリスト設定」「個別確認」「レビュー会議」「指摘整理」「是正・再確認」の箱を左から右へ並べる。下に「要求対応」「品質属性」「未解決事項」「リスク」を確認観点として並べる。日本語ラベル、講義ノート風。
\end{promptbox}

\subsection{品質属性}

設計品質に関わる属性には、モジュール性、保守容易性、移植性、テスト容易性、使いやすさ、正しさ、堅ろう性などがある。これらは、利害関係者の関心事の大きな部分を占める。

品質属性には、実行時に観察できるものと、実行時には直接見えにくいものがある。性能、セキュリティ、可用性、機能性、使いやすさは、実行時に観察しやすい。一方、変更容易性、移植性、再利用性、テスト容易性は、設計や保守の中で評価されることが多い。概念的一貫性、正しさ、完全性は、設計記述そのものを見て評価される。

\subsection{品質分析・評価技法}

設計品質を分析・評価するための技法には、静的分析、シミュレーション、プロトタイピング、レビュー、要求追跡などがある。

\begin{tabularx}{\linewidth}{L{36mm}Y}
\toprule
\textbf{技法} & \textbf{説明} \\
\midrule
設計レビュー・検査 & 設計記述を人が確認し、要求対応、未解決事項、設計上のリスクを見つける。 \\
シナリオに基づく評価 & 利用シナリオ、障害シナリオ、変更シナリオを使って、設計がどう反応するか確認する。 \\
要求追跡 & 要求が設計要素へ対応しているか、設計要素に根拠となる要求があるか確認する。 \\
静的分析 & 実行せずに設計やモデルを分析する。依存関係、循環、整合性、弱点を確認する。 \\
形式的分析 & 数学的モデルを使い、性質を厳密に検査する。安全性やセキュリティが重要な領域で有効である。 \\
シミュレーション & 設計モデルを使い、性能、タイミング、負荷、失敗時の挙動などを試す。 \\
プロトタイピング & 実現可能性、使いやすさ、技術的リスクを確認するために試作品を作る。 \\
\bottomrule
\end{tabularx}

\begin{todobox}
品質分析・評価技法の比較図を入れる。
\end{todobox}
\begin{promptbox}
白背景、モノクロ、A4本文幅の比較表風図を作成する。列は「レビュー」「静的分析」「形式的分析」「シミュレーション」「プロトタイプ」。行は「主な目的」「必要な成果物」「強み」「注意点」。各セルは短い日本語で埋める。講義ノート風、細線、印刷で読みやすい。
\end{promptbox}

\subsection{測度とメトリクス}

\term{測度}{Measure} と\term{メトリクス}{Metric} は、設計の大きさ、構造、複雑さ、品質を定量的に見るために使われる。測定対象は、設計方法によって異なる。

機能分解を中心とする設計では、構造チャートなどから、階層の深さ、モジュール数、呼び出し関係の複雑さを測る。オブジェクト指向設計では、クラス図などから、クラス数、継承の深さ、結合度、凝集度、メソッド数、メッセージ数などを測る。

測度は、設計の良し悪しを自動的に決めるものではない。測度は、注意すべき箇所を見つける手がかりである。たとえば、結合度が高い部品は変更の影響が大きい可能性がある。凝集度が低い部品は責務が散らばっている可能性がある。

\subsection{検証・妥当性確認・認証}

設計の分析と評価は、検証、妥当性確認、認証に関わる。

\begin{tabularx}{\linewidth}{L{32mm}Y}
\toprule
\textbf{活動} & \textbf{意味} \\
\midrule
検証 & 設計が、明示された要求を満たしているかを確認することである。 \\
妥当性確認 & その設計で作られるシステムが、顧客、利用者、運用者、保守担当者の期待を満たせるかを確認することである。 \\
認証 & 第三者が、設計が仕様や意図された利用に適合していると証明または表明することである。 \\
\bottomrule
\end{tabularx}

検証は「要求どおりか」を問う。妥当性確認は「本当に役に立つか」を問う。認証は「決められた基準に適合していると第三者が認めるか」を問う。安全性や規制が重要な領域では、設計段階から認証を見据えた証拠を残す必要がある。

\section{トピックと参考文献の対応}

この表は、SWEBOKが示す「トピック対参考資料の行列」を、学習用に簡略化したものである。詳しく学ぶときは、各トピックに対応する文献を読む。

\begin{longtable}{L{66mm}L{98mm}}
\toprule
\textbf{トピック} & \textbf{主な参考資料} \\
\midrule
\endhead
1. ソフトウェア設計の基礎 & Brooks, The Design of Design; Budgen, Software Design \\
1.1 設計思考 & Brooks; Budgen; Ross, Goodenough and Irvine \\
1.2 ソフトウェア設計の文脈 & Budgen; Sommerville \\
1.3 ソフトウェア設計の主要論点 & Brooks; Budgen; Sommerville; Bosch \\
1.4 ソフトウェア設計原則 & Budgen; Dijkstra; Parnas; ISO/IEC/IEEE 24765; IEEE Std 7000 \\
2. ソフトウェア設計プロセス & Budgen; Sommerville; ISO/IEC/IEEE 12207 \\
2.1 高水準設計 & Brooks; Budgen \\
2.2 詳細設計 & Budgen; Sommerville \\
3. ソフトウェア設計品質 & Budgen; Sommerville; Galster et al. \\
3.1 並行性 & Sommerville \\
3.2 制御とイベント処理 & Sommerville; Luckham; Mühl, Fiege and Pietzuch \\
3.3 データ永続化 & Sommerville \\
3.4 構成要素の分散 & Sommerville \\
3.5 誤り・例外処理と耐故障性 & Sommerville \\
3.6 統合と相互運用性 & Budgen \\
3.7 保証・セキュリティ・安全性 & Sommerville \\
3.8 変動性 & Galster et al.; Bosch \\
4. ソフトウェア設計の記録 & Booch, Rumbaugh and Jacobson; Budgen; Sommerville \\
4.1 モデルに基づく設計 & Budgen; Sommerville \\
4.2 構造設計記述 & Booch et al.; Budgen; Gamma et al.; Sommerville \\
4.3 振る舞い設計記述 & Booch et al.; Budgen; Gamma et al.; Sommerville \\
4.4 設計パターンとスタイル & Brooks; Budgen; Gamma et al.; Sommerville \\
4.5 専用言語と領域固有言語 & Sommerville; Kosar, Bohra and Mernik \\
4.6 設計根拠 & Brooks; Budgen; Sommerville; Parnas and Clements \\
5. ソフトウェア設計戦略と方法 & Sommerville; Budgen; Nielsen; Kiczales et al. \\
5.1 一般戦略 & Budgen \\
5.2 機能指向・構造化設計 & Budgen \\
5.3 データ中心設計 & Budgen; Sommerville \\
5.4 オブジェクト指向設計 & Budgen; Booch et al. \\
5.5 利用者中心設計 & Brooks; Nielsen \\
5.6 構成要素ベース設計 & Booch et al.; Budgen; Sommerville \\
5.7 イベント駆動設計 & Sommerville; Luckham; Mühl et al. \\
5.8 アスペクト指向設計 & Booch et al.; Sommerville; Kiczales et al. \\
5.9 制約に基づく設計 & Brooks \\
5.10 ドメイン駆動設計 & Budgen \\
5.11 その他の方法 & Sommerville \\
6. 設計品質分析と評価 & Budgen; Sommerville; Parnas and Weiss \\
6.1 設計レビューと監査 & Budgen; Parnas and Weiss \\
6.2 品質属性 & Sommerville \\
6.3 品質分析・評価技法 & Sommerville; Software Quality KA \\
6.4 測度とメトリクス & Budgen; Sommerville \\
6.5 検証・妥当性確認・認証 & Sommerville; Software Quality KA \\
\bottomrule
\end{longtable}

\section{発展的な読み物}

SWEBOK 03章で紹介される発展的な読み物を、日本語で読む目的に合わせて整理する。

\begin{tabularx}{\linewidth}{L{52mm}Y}
\toprule
\textbf{文献} & \textbf{読むと得られること} \\
\midrule
Brooks, \textit{The Design of Design} & ソフトウェア工学の先駆者による設計論である。設計活動、設計判断、事例を広く学べる。 \\
Budgen, \textit{Software Design} & ソフトウェア設計を体系的に学ぶための中心的な教科書である。構造、振る舞い、設計方法、評価を幅広く扱う。 \\
Gamma et al., \textit{Design Patterns} & 再利用可能なオブジェクト指向設計パターンの古典である。設計語彙を学ぶのに有効である。 \\
Booch, Rumbaugh and Jacobson, \textit{The Unified Modeling Language User Guide} & UMLを使った構造・振る舞いの表現を学べる。設計記述の記法を理解するのに役立つ。 \\
Sommerville, \textit{Software Engineering} & 設計だけでなく、要求、実装、テスト、品質、分散システムなどを含むソフトウェア工学全体を学べる。 \\
Parnas, “On the Criteria To Be Used In Decomposing Systems Into Modules” & モジュール分解と情報隠蔽の古典的論文である。設計原則の背景を学べる。 \\
IEEE Std 7000-2021 & 倫理的関心をシステム設計で扱うための標準である。AIや自律システムの設計で参考になる。 \\
\bottomrule
\end{tabularx}

\section{参考文献}

\begin{enumerate}
  \item G. Booch, J. Rumbaugh, and I. Jacobson, \textit{The Unified Modeling Language User Guide}, 2nd edition, Addison-Wesley, 2005.
  \item J. Bosch, \textit{Design and Use of Software Architectures: Adopting and Evolving a Product-Line Approach}, ACM Press, 2000.
  \item F. Brooks, \textit{The Design of Design}, Addison-Wesley, 2010.
  \item D. Budgen, \textit{Software Design: Creating Solutions for Ill-Structured Problems}, 3rd edition, CRC Press, 2021.
  \item E. W. Dijkstra, “On the Role of Scientific Thought,” 1974.
  \item M. Galster, D. Weyns, D. Tofan, B. Michalik, and P. Avgeriou, “Variability in Software Systems - A Systematic Literature Review,” \textit{IEEE Transactions on Software Engineering}, 40(3), 2014.
  \item E. Gamma et al., \textit{Design Patterns: Elements of Reusable Object-Oriented Software}, Addison-Wesley, 1994.
  \item S. Gregor and D. Jones, “The Anatomy of a Design Theory,” Association for Information Systems, 2007.
  \item IEEE Std 7000-2021, \textit{IEEE Standard Model Process for Addressing Ethical Concerns during System Design}.
  \item ISO/IEC/IEEE 12207, \textit{Systems and Software Engineering - Software Life Cycle Processes}.
  \item ISO/IEC/IEEE 24765:2017, \textit{Systems and Software Engineering - Vocabulary}, 2nd ed., 2017.
  \item G. Kiczales et al., “Aspect-Oriented Programming,” Proc. 11th European Conf. Object-Oriented Programming, Springer, 1997.
  \item T. Kosar, S. Bohra, and M. Mernik, “Domain-Specific Languages: A Systematic Mapping Study,” \textit{Information and Software Technology}, 71, 77-91, 2016.
  \item D. Luckham, \textit{The Power of Events: an Introduction to Complex Event Processing}, Addison-Wesley, 2002.
  \item G. Mühl, L. Fiege, and P. Pietzuch, \textit{Distributed Event-Based Systems}, Springer-Verlag, 2006.
  \item J. Nielsen, \textit{Usability Engineering}, Morgan Kaufmann, 1994.
  \item D. L. Parnas, “On the Criteria To Be Used In Decomposing Systems Into Modules,” \textit{Communications of the ACM}, 15(12), 1053-1058, 1972.
  \item D. L. Parnas and P. C. Clements, “A Rational Design Process: How and Why to fake it,” \textit{IEEE Transactions on Software Engineering}, 12(2), 251-257, 1986.
  \item D. L. Parnas and D. M. Weiss, “Active Design Reviews: Principles and Practices,” \textit{Journal of Systems and Software}, 7, 259-265, 1987.
  \item D. T. Ross, J. B. Goodenough, and A. Irvine, “Software Engineering: Process, Principles, and Goals,” \textit{IEEE Computer}, May 1975.
  \item I. Sommerville, \textit{Software Engineering}, 10th edition, Pearson, 2016.
\end{enumerate}

\section{画像 TODO 一覧}

本文に配置した画像 TODO を再掲する。実際に図を作る場合は、各TODOのプロンプトをそのまま画像生成に使うと、A4資料に合う図を作りやすい。

\begin{enumerate}
  \item 03章の全体地図。
  \item 設計思考の五段階図。
  \item 設計と他知識領域の関係図。
  \item 主要な設計原則の関係図。
  \item 設計の三段階図。
  \item 設計品質の関心事マップ。
  \item 文書中心設計とモデル中心設計の比較図。
  \item 構造設計記述と振る舞い設計記述の分類図。
  \item 設計パターン分類図。
  \item 設計戦略と中心概念の対応図。
  \item 設計レビューの流れ。
  \item 品質分析・評価技法の比較図。
\end{enumerate}

\section{章末確認問題}

\begin{enumerate}
  \item ソフトウェア設計という語が持つ四つの意味を説明せよ。
  \item 設計思考の五段階を、自分の言葉で説明せよ。
  \item 高水準設計と詳細設計の違いを、外部向き・内部向きという観点から説明せよ。
  \item 抽象化、モジュール化、カプセル化の違いを説明せよ。
  \item 結合度を低くし、凝集度を高くすることが望ましい理由を説明せよ。
  \item 並行性を設計するとき、同期や排他制御が必要になる理由を説明せよ。
  \item 構造設計記述と振る舞い設計記述の例を、それぞれ三つ挙げよ。
  \item 設計パターンが「語彙」として役立つ理由を説明せよ。
  \item ドメイン駆動設計で、共有語彙が重要になる理由を説明せよ。
  \item 設計レビュー、静的分析、プロトタイピングはどのように使い分けるか。
\end{enumerate}

\begin{pointbox}{解答の方向性}
1は「分野、プロセス、結果、ライフサイクル段階」である。3は、高水準設計が外部とのやり取りや主要構成要素の責務を扱い、詳細設計が内部モジュール、データ構造、アルゴリズムを扱う点を説明する。5は、変更影響の局所化と理解容易性を軸に説明できる。7は、構造ならクラス図、コンポーネント図、ERD、配置図など、振る舞いなら活動図、シーケンス図、DFD、状態遷移図、決定表などを挙げればよい。
\end{pointbox}

\end{document}
